\documentclass[10pt,a4paper]{amsart}
\usepackage[margin=19mm]{geometry}
\usepackage{amsmath,amssymb,mathtools}
\usepackage[scr=rsfs,cal=euler]{mathalfa}
\usepackage{xfrac}
\usepackage[unicode,hidelinks,bookmarks=false]{hyperref}
\newcommand{\ie}{i.e.\ }
\newcommand{\cf}{cf.\ }
\newcommand{\resp}{respectively\ }
\newcommand{\Subsection}{\subsection}
\providecommand{\nobreakdash}{\nobreak\hbox{-}}
\setlength{\emergencystretch}{2em}
\multlinegap0pt
\usepackage{amssymb}
\newtheorem{theorem}{Theorem}
\title{From Nonsmooth Minima to Smooth Branches via Heat Kernel Regularization}
\author{Hyeontae Jo}
\date{Source: arXiv:2604.05372v1 (CC BY 4.0)}
\begin{document}
\maketitle

\noindent\emph{Source and licence.} From Nonsmooth Minima to Smooth Branches via Heat Kernel Regularization. Version arXiv:2604.05372v1 (CC BY 4.0), distributed by arXiv under the Creative Commons Attribution 4.0 International licence, http://creativecommons.org/licenses/by/4.0/, as recorded in the arXiv licence metadata for this exact version and shown on the abstract page https://arxiv.org/abs/2604.05372v1. Full licence notice: \texttt{licenses/arxiv-2604.05372-CC-BY-4.0.txt}.\par\smallskip

\noindent\textit{Editorial scope.} Counted unit: the article's theorem thm:continuation-alternative (source0.tex lines 614--632). The article's complete original proof of it is delivered with it (source0.tex lines 634--708, one \texttt{proof} environment of 75 lines). No mathematics is written, repaired or re-derived: the statement and the proof are the article's own lines, in the article's own notation, and no retained line is substituted or re-flowed.  No further source-local context is needed: every cross-reference of every retained line resolves inside the two delivered spans.  Nothing was omitted from any retained span, so the package carries no omission record.  No retained line contains a citation, so the wrapper carries no bibliography and no bibliography entry.  No retained line contains a figure environment, an image-inclusion command or drawing code, so no image is delivered and the package declares no asset dependency.  Editorial formatting only, in addition: an AMS \texttt{amsart} preamble that restates the article's own declarations and macros used by the retained lines, namely the environments \texttt{theorem} on separate counters, together with the standard packages those lines need.  The retained environments print in this document as Theorem 1 in order of appearance, whereas the article prints the same items with the numbering of its own full document.  The complete native source of the article as distributed in the arXiv e-print for this exact version is bundled unchanged as \texttt{sources/197951/source0.tex}.
\begin{theorem}\label{thm:continuation-alternative}
Let $x:(0,T_{\max})\to\mathbb R^n$ be a nondegenerate local minimizing branch with $0<T_{\max}<\infty$, and assume that
$$
x_t\to \bar x
\qquad\text{as }t\to T_{\max}^-
$$
for some $\bar x\in\mathbb R^n$. Then the following are equivalent:
\begin{enumerate}
\item[(i)] The branch $x_t$ extends across $T_{\max}$ as a nondegenerate local minimizing branch.
\item[(ii)] The Hessian remains uniformly nondegenerate near $T_{\max}$, namely,
$$
\liminf_{t\to T_{\max}^-}\lambda_{\min}\bigl(\nabla_x^2(P_t f)(x_t)\bigr)>0.
$$
\end{enumerate}
In particular, if the branch is maximal, then
$$
\liminf_{t\to T_{\max}^-}\lambda_{\min}\bigl(\nabla_x^2(P_t f)(x_t)\bigr)=0.
$$
\end{theorem}

\begin{proof}
Assume first that (ii) holds. Then there exist $c_0>0$ and $t_1<T_{\max}$ such that
$$
\lambda_{\min}\bigl(\nabla_x^2(P_t f)(x_t)\bigr)\ge c_0
\qquad\text{for all }t\in(t_1,T_{\max}).
$$
Set
$$
F(t,x):=\nabla_x(P_t f)(x).
$$
Since $P_t f$ is smooth in $(t,x)$ for $t>0$, the map $F$ is continuous in $t$ and $C^1$ in $x$ near $(T_{\max},\bar x)$. Because $x_t\to\bar x$ and $F(t,x_t)=0$ for all $t<T_{\max}$, continuity gives
$$
F(T_{\max},\bar x)=0.
$$
Moreover, by continuity of the Hessian and the lower bound along the branch,
$$
\nabla_xF(T_{\max},\bar x)=\nabla_x^2(P_{T_{\max}}f)(\bar x)\succ 0.
$$
Hence $\nabla_xF(T_{\max},\bar x)$ is invertible. By the implicit function theorem, there exist $\varepsilon>0$, a neighborhood $U$ of $\bar x$, and a unique $C^1$ map
$$
y:(T_{\max}-\varepsilon,T_{\max}+\varepsilon)\to U
$$
such that
$$
F(t,y(t))=0
\qquad\text{and}\qquad
y(T_{\max})=\bar x.
$$
Since $x_t\to\bar x$, we have $x_t\in U$ for all $t$ sufficiently close to $T_{\max}$, and by uniqueness in the implicit function theorem,
$$
x_t=y(t)
\qquad\text{for all }t\in(T_{\max}-\varepsilon,T_{\max}).
$$
Shrinking $\varepsilon$ if necessary, continuity of the Hessian implies
$$
\nabla_x^2(P_t f)(y(t))\succ 0
\qquad\text{for all }t\in(T_{\max}-\varepsilon,T_{\max}+\varepsilon).
$$
Thus $y$ extends the original branch across $T_{\max}$ as a nondegenerate local minimizing branch. Therefore (ii)$\Rightarrow$(i).

Conversely, assume that (i) holds. Then there exist $\varepsilon>0$ and a nondegenerate local minimizing branch
$$
\widetilde x:(T_{\max}-\varepsilon,T_{\max}+\varepsilon)\to\mathbb R^n
$$
such that
$$
\widetilde x_t=x_t
\qquad\text{for all }t\in(T_{\max}-\varepsilon,T_{\max}).
$$
Since
$$
\nabla_x^2(P_t f)(\widetilde x_t)\succ 0
\qquad\text{for all }t\in(T_{\max}-\varepsilon,T_{\max}+\varepsilon),
$$
continuity of the Hessian implies that there exist $c_0>0$ and $\delta>0$ such that
$$
\lambda_{\min}\bigl(\nabla_x^2(P_t f)(\widetilde x_t)\bigr)\ge c_0
\qquad\text{for all }t\in(T_{\max}-\delta,T_{\max}+\delta).
$$
Hence
$$
\liminf_{t\to T_{\max}^-}\lambda_{\min}\bigl(\nabla_x^2(P_t f)(x_t)\bigr)\ge c_0>0.
$$
Thus (i)$\Rightarrow$(ii).

Finally, if the branch is maximal, then (i) fails. Therefore (ii) also fails. Since
$$
\lambda_{\min}\bigl(\nabla_x^2(P_t f)(x_t)\bigr)>0
\qquad\text{for all }t\in(0,T_{\max}),
$$
we necessarily have
$$
\liminf_{t\to T_{\max}^-}\lambda_{\min}\bigl(\nabla_x^2(P_t f)(x_t)\bigr)=0.
$$
\end{proof}

\end{document}
